\subsection{Comparison with simulated annealing}
\label{subsec:sa_comparison}
We compare two symmetric-TT variants with two simulated annealing (SA)
baselines on assignment and the low-rank strip Set Partitioning family.
Each family and size $n\in\{10,12,14\}$ contains 40 paired instances with
the same 400 initial feasible draws. Separate runs use 10, 50 and 200-second
budgets on one CPU core per solver; both quadratic and entropic objectives
are evaluated as black boxes. We show $n=14$, 200-second trajectories in
\cref{fig:sa_time,fig:sa_calls}; the corresponding $n=10,12$ curves and
complete protocol appear in \cref{app:sa_protocol}.

\textbf{Information and model variants.} Data-initialized EELS builds a
symmetric TT from the current selected data and expands its graph with
$\gamma=1$. Full-indicator TT starts from the exact indicator of all feasible
solutions. In this experiment EELS rebuilds its model, and full-indicator TT
resets its amplitudes, at every outer iteration; this differs from the
persistent-amplitude EELS--Best Cost comparison in
\cref{subsec:eels_optimization}. Informed SA receives explicit
feasibility-preserving moves. Penalty SA receives no such moves: it flips a
random bit and minimizes $C(x)+\alpha\|Ax-b\|_2^2$ with a fixed scale-based
choice of $\alpha$. EELS and penalty SA both use the constraints $A,b$ and
the initial data, but exploit them differently: EELS enforces feasibility
structurally, whereas penalty SA explores infeasible points. This is a closer
information comparison than informed SA, not an assertion of identical
algorithmic prior knowledge. Full-indicator construction likewise uses
additional problem structure.

\textbf{Finite-time quality.} EELS outperforms the tested penalty-SA
configuration in every 200-second group in mean paired final quality;
all twelve one-sided comparisons remain significant after Holm correction
over the 72 TT-versus-penalty comparisons across sizes, objectives and budgets
(\cref{tab:sa_penalty_tests}). Penalty SA does improve its initial feasible
best in 630 of 1,440 runs, including improvements in all twelve 200-second
groups, but is substantially less effective within these budgets. This
conclusion concerns the specified penalty, single-bit proposal and cooling
schedule; it does not establish superiority over optimally tuned or adaptive
penalty methods.

The EELS curves typically gain quickly and then plateau, whereas penalty SA
continues to make slower gains on the displayed horizon. This is consistent
with EELS acting as a rapid intensification step on data-defined support.
A flat incumbent curve does not prove a local optimum, and an improving
mean does not establish that individual SA chains never stagnate. A possible
hybrid would supply feasible solutions from SA to a subsequent EELS stage.
Such complementarity is motivated by the SA-seeded TN-GEO booster studied
by \citet{Alcazar2024}, but the proposed SA--EELS hybrid is not tested here.

\textbf{Rank and informed SA.} On assignment, informed SA has better final
cost on average in every group, with no individual TT wins (some ties).
Its transpositions encode
the permutation structure directly, while EELS receives no feasible-move
oracle. Full-indicator TT approaches informed SA at $n=10$, but loses more
at larger sizes. Its exact maximum rank grows from 461 to 1,715 to 6,434;
construction and training overhead are consistent with reduced search
throughput. The present comparison does not isolate rank from size,
objective landscape or reset behavior, so it is not a causal timing study.
On low-rank Set Partitioning, the TT endpoints are much closer to informed
SA, particularly for the entropic objective, with many paired ties and
some TT wins. Neither TT variant has a Holm-significant advantage over
informed SA; failure to reject a difference is not an equivalence result.
Thus the observed rank contrast explains a plausible computational mechanism,
not a universal ordering of optimizers.

\textbf{Objective-query efficiency.} The evaluation-axis curves tell a
different story from wall time. On Set Partitioning the TT methods commonly
reach comparable mean gains with far fewer objective calls; for example,
in the $n=14$ entropic 200-second runs, informed SA averages 58.3 million
calls versus 0.66 million for EELS and 0.54 million for full-indicator TT.
Both TT variants also have higher mean gains than informed SA at the
100,000-call slice on every Set Partitioning group of the 200-second runs.
Full-indicator TT has this advantage on assignment at $n=10,12$, while
EELS does not. These are descriptive slices of time-limited trajectories:
SA cooling depends on elapsed time, and every curve is truncated to a query
range reached by all 40 tasks. They neither constitute a separately tuned
equal-query experiment nor guarantee a wall-time advantage for expensive
objectives. All calls, including repeats and infeasible penalty-SA queries,
are counted; no solver receives incremental objective updates.

\begin{figure*}[t]
\centering
\includegraphics[width=0.96\textwidth]{pictures/SA_n14_time.pdf}
\caption{Mean feasible-objective gain versus optimization time at $n=14$
($N=196$), with 40 paired tasks per panel and a 200-second nominal budget.
Top: assignment; bottom: spatial Set Partitioning. Left: quadratic cost;
right: entropic risk. Gain is $100(C_0-C_{\mathrm{best}})/|C_0|$ per task,
then averaged. Shading gives pointwise asymmetric 95\% bootstrap intervals
for the mean, not task standard deviations. The same initial data are used
by all four methods; penalized or infeasible costs never enter the ordinate.
EELS and full-indicator TT use the reset/rebuild protocol described in
\cref{app:sa_protocol}, not persistent-amplitude training.}
\label{fig:sa_time}
\end{figure*}

\begin{figure*}[t]
\centering
\includegraphics[width=0.96\textwidth]{pictures/SA_n14_calls.pdf}
\caption{The same $n=14$, 200-second runs as \cref{fig:sa_time}, plotted
against cumulative objective calls (logarithmic axis). Calls include the
initial evaluations, duplicates, and infeasible penalty-SA points.
Each method stops at the smallest final call count among its 40 tasks
(endpoint marker); no stopped trajectory is extrapolated. Comparisons
are meaningful only on overlapping query ranges. Confidence intervals and
gain normalization match \cref{fig:sa_time}. These curves describe
query efficiency under time-based cooling, rather than a fixed-query
benchmark.}
\label{fig:sa_calls}
\end{figure*}

